% !TEX program = xelatex
% ============================================================
%  دعوة زفاف — Wedding Invitation Template
%  قالب دعوة زفاف أنيق بإطار زخرفي وعناصر تزيينية.
%  Compile with: xelatex main.tex (run twice)
% ============================================================
%  Features:
%    - Decorative TikZ border with corner flourishes
%    - "بسم الله الرحمن الرحيم" header
%    - Bride and groom names in large Amiri font
%    - Date, time, venue and RSVP contact fields
%    - Single page with generous whitespace
% ============================================================

\documentclass[11pt,a4paper]{article}

% ---- Geometry ----
\usepackage[a4paper,margin=2.5cm,top=2.5cm,bottom=2.5cm,headheight=16pt]{geometry}

% ---- Core packages (order matters for bidi) ----
\usepackage{graphicx}
\usepackage{array}
\usepackage{booktabs}
\usepackage{tabularx}
\usepackage{enumitem}
\usepackage{setspace}
\usepackage{xcolor}
\usepackage{tikz}
\usetikzlibrary{calc,decorations.markings,decorations.pathmorphing}

% ---- RTL / Arabic language support (must load before hyperref) ----
\usepackage{bidi}
\usepackage[no-math]{fontspec}
\usepackage[quiet]{polyglossia}

% ---- Fonts ----
\setmainfont[Scale=1]{NotoKufiArabic-Regular.ttf}
\newfontfamily\arabicfont[Script=Arabic,Scale=0.95]{NotoKufiArabic-Regular.ttf}
\newfontfamily\englishfont[Scale=0.95]{TeX Gyre Termes}

% ---- Languages ----
\setdefaultlanguage[calendar=gregorian,hijricorrection=1]{arabic}
\setotherlanguage[variant=british]{english}

% ---- Hyperref (load after polyglossia/bidi) ----
\usepackage[breaklinks,hidelinks]{hyperref}
\hypersetup{
  pdftitle={دعوة زفاف},
  pdfauthor={الاسم}
}

% ---- Colors ----
\definecolor{invgold}{HTML}{8B6F2A}
\definecolor{invcream}{HTML}{FBF7EF}
\definecolor{invink}{HTML}{2B2B2B}

% ---- Line spacing ----
\setstretch{1.4}

\pagestyle{empty}

% ---- Decorative corner flourish ----
\newcommand{\cornerflourish}[2]{%
  \begin{scope}[shift={(#1,#2)}]
    \draw[invgold,line width=1pt] (0,0) -- (1.6,0);
    \draw[invgold,line width=1pt] (0,0) -- (0,1.6);
    \draw[invgold,line width=0.6pt] (0.25,0.25) -- (1.2,0.25) -- (1.2,1.2);
    \draw[invgold,line width=0.6pt] (0.25,0.25) -- (0.25,1.2) -- (1.2,1.2);
    \fill[invgold] (0,0) circle (2pt);
    \fill[invgold] (1.6,0) circle (1.5pt);
    \fill[invgold] (0,1.6) circle (1.5pt);
  \end{scope}%
}

\begin{document}
\color{invink}

% ---- Decorative border frame ----
\begin{tikzpicture}[remember picture,overlay]
  \draw[invgold,line width=1.2pt]
    ($(current page.north west)+(1.2cm,-1.2cm)$)
    rectangle
    ($(current page.south east)+(-1.2cm,1.2cm)$);
  \draw[invgold,line width=0.5pt]
    ($(current page.north west)+(1.5cm,-1.5cm)$)
    rectangle
    ($(current page.south east)+(-1.5cm,1.5cm)$);
  % Corner flourishes (mirror for each corner)
  \begin{scope}[shift={(current page.north west)}]
    \cornerflourish{1.2}{-1.2}
  \end{scope}
  \begin{scope}[shift={(current page.north east)},xscale=-1]
    \cornerflourish{1.2}{-1.2}
  \end{scope}
  \begin{scope}[shift={(current page.south west)},yscale=-1]
    \cornerflourish{1.2}{-1.2}
  \end{scope}
  \begin{scope}[shift={(current page.south east)},xscale=-1,yscale=-1]
    \cornerflourish{1.2}{-1.2}
  \end{scope}
\end{tikzpicture}

\vspace*{0.4cm}

\begin{center}

% ---- Bismillah header ----
{\fontsize{24pt}{30pt}\selectfont\color{invgold}
بسم الله الرحمن الرحيم\par}

\vspace{0.7cm}

{\fontsize{14pt}{18pt}\selectfont
\textarabic{وَمِنْ آيَاتِهِ أَنْ خَلَقَ لَكُم مِّنْ أَنفُسِكُمْ أَزْوَاجًا لِّتَسْكُنُوا إِلَيْهَا}\par}
{\fontsize{12pt}{16pt}\selectfont\color{invgold}
\textarabic{وَجَعَلَ بَيْنَكُم مَّوَدَّةً وَرَحْمَةً}\par}

\vspace{0.9cm}

% ---- Invitation line ----
{\fontsize{18pt}{24pt}\selectfont
يسرّنا دعوتكم لحضور حفل زفاف\par}

\vspace{0.7cm}

% ---- Bride and groom names ----
% TODO: استبدل الأسماء بأسماء العروسين
{\fontsize{26pt}{32pt}\selectfont\color{invgold}
اسم العروس\par}
\vspace{0.1cm}
{\fontsize{16pt}{20pt}\selectfont
\&\par}
\vspace{0.1cm}
{\fontsize{26pt}{32pt}\selectfont\color{invgold}
اسم العريس\par}

\vspace{0.8cm}

% ---- Decorative divider ----
\begin{tikzpicture}
  \draw[invgold,line width=0.8pt] (-4,0) -- (-0.6,0);
  \draw[invgold,line width=0.8pt] (0.6,0) -- (4,0);
  \fill[invgold] (0,0) circle (3pt);
  \draw[invgold,line width=0.5pt] (0,0) circle (7pt);
\end{tikzpicture}

\vspace{0.6cm}

% ---- Event details ----
% TODO: املأ تاريخ ووقت ومكان الحفل
{\fontsize{13pt}{18pt}\selectfont
\begin{tabular}{r@{\hspace{1.2cm}}l}
\textbf{التاريخ:} & يوم الجمعة، ١٥ شعبان ١٤٤٦هـ \\
\textbf{الوقت:} & الساعة السابعة مساءً \\
\textbf{المكان:} & قاعة المناسبات — المدينة \\
\end{tabular}\par}

\vspace{0.8cm}

% ---- RSVP ----
{\fontsize{12pt}{16pt}\selectfont
\textbf{يرجى تأكيد الحضور (RSVP):}\par}
\vspace{0.1cm}
{\fontsize{11pt}{15pt}\selectfont\color{invgold}
٠٥٥٥٥٥٥٥٥٥ — الاسم\par}

\end{center}

\end{document}
